\begin{longtable}{|p{1.5cm}|p{5cm}|p{7cm}|}
\caption{Kebutuhan Non-Fungsional}
\label{tbl:kebutuhan-nonfungsional} \\
\hline
\textbf{ID} & \textbf{Kebutuhan} & \textbf{Penjelasan} \\
\hline
\endfirsthead

\hline
\textbf{ID} & \textbf{Kebutuhan} & \textbf{Penjelasan} \\
\hline
\endhead

\hline
\endfoot

NFR-01 & Kinerja &
Sistem harus mampu melakukan \textit{crawling} dan eksfiltrasi data tanpa menimbulkan konsumsi
\textit{resource} (CPU, memori, jaringan) yang tinggi agar tidak memicu alarm \textit{anti-malware}
atau monitoring forensik. \\
\hline

NFR-02 & Reliabilitas &
Segala \textit{failure} harus ditangani secara \textit{silent} (tidak menimbulkan pesan
\textit{error}/\textit{alert} pada \textit{endpoint} atau sistem keamanan). \\
\hline

NFR-03 & Keamanan lingkungan uji &
Sistem harus memastikan seluruh eksekusi \textit{malware} berlangsung dalam lingkungan yang benar-benar terisolasi dan tidak terhubung dengan sistem produksi. \\
\hline

NFR-04 & Kepatuhan terhadap regulasi dan etika penelitian &
Sistem wajib mematuhi regulasi keamanan siber, etika penelitian, serta tidak menggunakan
\textit{malware} untuk tujuan di luar pengujian akademik terisolasi. \\
\hline

NFR-05 & Isolasi Dampak &
Efek dan aktivitas \textit{malware} harus benar-benar terbatas pada lingkungan uji,
tidak boleh berdampak pada sistem luar sehingga integritas lab tetap terjaga. \\
\hline

\end{longtable}